\subNsection{1}{Zusammenfassung der Projektergebnisse}{Alle}

Kristian Popov: Mein Schwerpunkt im Projekt lag auf den Themen
\emph{Endliche Automaten} und \emph{Reguläre Sprachen}. Für Endliche Automaten
habe ich eine interaktive Umgebung umgesetzt, in der deterministische und
nichtdeterministische Automaten erstellt und Eingabewörter schrittweise ausgeführt
werden können. Übergangstabelle und Zustandsdiagramm zeigen dabei denselben Ablauf,
sodass die formale Definition und ihre grafische Darstellung miteinander verbunden
werden. Ein Stapeltest prüft mehrere Wörter auf einmal und macht dadurch sichtbar, ob
ein Automat die beabsichtigte Sprache erkennt und nicht nur das gerade ausprobierte
Wort. Ergänzt wird das Thema durch Beispiele, Übungen mit Rückmeldung und die
Umwandlung rechtslinearer Grammatiken in Automaten.

Im Thema Reguläre Sprachen stehen drei Arbeitsweisen zur Verfügung: Studierende können
die von einer Grammatik erzeugten Wörter aufzählen, aus einer Sprachbeschreibung eine
passende Grammatik erzeugen und einzelne Wörter auf ihre Zugehörigkeit zur Sprache
prüfen. Bei der Wortprüfung wird neben dem Ergebnis auch dessen Grund angezeigt.
Übungen ermöglichen auch hier eine direkte Überprüfung der eigenen Lösung. In beiden
Themen lassen sich Automaten, Grammatiken und Sprachbeschreibungen über die gemeinsame
Import- und Exportfunktion der Plattform weitergeben. Bei beiden Themen war mir wichtig,
die Fachlogik klar von der Benutzeroberfläche zu trennen. Dadurch konnten gemeinsame
Grammatik- und Automatenfunktionen wiederverwendet werden, während die einzelnen
Funktionen testbar und übersichtlich blieben.

Duy Tan Pham: Mein fachlicher Schwerpunkt lag auf dem Themenbereich
\emph{Disjunktive Normalform}. Dort können Studierende boolesche Ausdrücke und
Wahrheitstabellen untersuchen, die kanonische DNF schrittweise herleiten und Ausdrücke
vereinfachen. Beispiele zeigen einen vollständigen Lösungsweg. In den Übungen bearbeiten
die Lernenden dagegen Wahrheitstabellen oder DNF-Ausdrücke selbst und erhalten eine
fachlich berechnete Rückmeldung. Dabei werden auch logisch gleichwertige Schreibweisen
erkannt.

Die Minimierung liefert eine DNF mit minimaler Literalzahl. Zusätzlich kann für bis zu
vier Variablen eine global kleinste Formel aus UND- und ODER-Verknüpfungen bestimmt
werden. Die Fachlogik ist von der Benutzeroberfläche getrennt und durch automatisierte
Tests abgesichert. Der DNF-Bereich nutzt außerdem den gemeinsamen Themenaufbau sowie die
vorhandenen Funktionen für Import und Export. Dadurch fügt er sich in Bedienung und
Darstellung in die gesamte Lernplattform ein.
